%==============================================================================
% MUWS 2026 — Cyber/endpoint-security reframe of the combined paper.
% Same experiments, same numbers; framing: offline on-device harmful media scanner.
% Compile on Overleaf: upload this file + refs.bib; compiler = pdfLaTeX; run twice.
% NOTE: ~6-7 pages. See trimming notes at the bottom if CFP enforces 4 pages.
%==============================================================================
\documentclass[sigconf,nonacm]{acmart}
\settopmatter{printacmref=false}
\usepackage{booktabs}
\usepackage{graphicx}

\begin{document}

\title{No Cloud Required:\\
On-Device Small Vision--Language Models for Offline\\
Detection of Harmful Multimodal Content}

\author{Anonymous Author(s)}
\affiliation{\institution{Anonymous Institution}\country{}}

\begin{abstract}
Content moderation systems that depend on cloud APIs offer no protection once a device
goes offline---a growing gap exploited by harmful media distributed via physical storage,
peer-to-peer channels, or simply downloaded for later offline viewing. Parental-control
apps, corporate endpoint-security tools, and school-device management all require a scanner
that can run \emph{entirely on the device} without a network connection. Compact
vision--language models (VLMs, $\leq$2B parameters) are a natural candidate---but their
behaviour as offline content scanners is poorly understood. We evaluate three small VLMs
(SmolVLM-256M, SmolVLM-500M, Gemma-4-E2B) on two harmful-content benchmarks (Hateful Memes
and MMSD2.0) under zero-shot and parameter-efficient (LoRA) conditions, and report four
findings with direct implications for on-device deployment. \textbf{(1) Memory budget, not
raw parameter count, governs feasibility:} the 0.5B SmolVLM \emph{Pareto-dominates} the
$\sim$4$\times$-larger Gemma-4-E2B---higher detection quality at $\sim$1/7 the VRAM.
\textbf{(2) Zero-shot models fail silently:} they emit a near-constant label (SmolVLM
floods ``yes'', 86--98\%; Gemma floods ``no'', 100\%), making them functionally equivalent
to a broken scanner---and standard metrics \emph{hide} this failure.
\textbf{(3) A cheap policy adapter ($<$2\% extra parameters) fixes the collapse first, then
improves detection---but the gain is 7$\times$ larger for sarcastic content than for
adversarially designed hate memes.} \textbf{(4) A hate adapter transfers to sarcasm
detection; the reverse does not hold}---useful guidance for bootstrapping a new content
category from an existing one. We release a reproducible evaluation harness and distil
concrete deployment recommendations.
\end{abstract}

\keywords{on-device content moderation, offline scanning, harmful memes, multimodal safety,
endpoint security, vision--language models, LoRA, cross-task transfer}

\maketitle

%------------------------------------------------------------------------------
\section{Introduction}
%------------------------------------------------------------------------------
Harmful multimodal content---hateful memes, inappropriate images with overlaid text,
ironic abuse---is increasingly distributed \emph{outside} the reach of server-side
moderation. A child may download videos and image sets for offline viewing before a device
is handed to them; harmful memes spread via USB drives, AirDrop, or encrypted peer-to-peer
channels that bypass network inspection; an air-gapped corporate workstation can accumulate
policy-violating media with no cloud service ever observing it. In all these scenarios,
the only viable moderation path is a scanner that runs locally on the endpoint itself.

Existing endpoint content-moderation tools rely primarily on perceptual hashing of known
contraband (e.g.\ PhotoDNA for CSAM), which offers no coverage for \emph{novel} or
\emph{borderline} harmful material that has not yet been catalogued. A semantic,
generalisation-capable model is needed. Small VLMs are the only class of multimodal models
that can plausibly run inference on consumer-grade hardware (a laptop or tablet GPU or
even CPU) within the memory and latency budgets of a background scanning process. Yet their
behaviour at the $\leq$2B scale on harmful-content detection tasks is largely uncharted.

This paper characterises the practical viability of small VLMs as on-device harmful-content
scanners via four interlocking findings. The picture that emerges is sobering but
actionable: zero-shot models are \emph{not} deployable as scanners---they silently fail in
a way that would give a false sense of security---but a single round of lightweight
adaptation (LoRA, $<$2\% extra parameters, negligible inference overhead) converts the best
small model into a functional classifier. The catch is that this works well for some
content categories and poorly for others, and the gap is explainable.

\noindent\textbf{Contributions.}
\begin{itemize}
\item A memory-footprint Pareto analysis showing a 0.5B model Pareto-dominates a 2B one
for on-device deployment (\S\ref{sec:scale}).
\item Evidence of \emph{label-prior collapse}---a silent failure mode where a VLM acts as
a broken scanner---and a demonstration that standard metrics hide it
(\S\ref{sec:collapse}).
\item A matched-budget LoRA study showing that policy adaptation first \emph{fixes} the
broken scanner, then delivers a 7$\times$ difficulty gap across content categories
(\S\ref{sec:adapt}).
\item A cross-category transfer matrix showing that a hate adapter bootstraps sarcasm
detection but not vice versa (\S\ref{sec:transfer}).
\end{itemize}

%------------------------------------------------------------------------------
\section{Related Work}
%------------------------------------------------------------------------------
\textbf{Harmful-content benchmarks.} The Hateful Memes Challenge~\cite{kiela2020hateful}
introduced ``benign confounders'' so that unimodal or shallow models fail; it closely
resembles content that circulates in the wild because images and text are independently
benign but jointly toxic. MMSD2.0~\cite{qin2023mmsd2} re-annotates the multimodal sarcasm
corpus of \citet{cai2019multimodal} to remove spurious text-only cues, making it resistant
to text-alone classifiers. HarMeme~\cite{pramanick2021harmeme} extends hate memes to new
political and social domains. Together these benchmarks represent the kinds of novel,
semantically grounded harmful images that hash-based endpoint tools cannot catch.

\textbf{Efficient VLMs for edge deployment.} SmolVLM~\cite{marafioti2025smolvlm}, built on
Idefics3~\cite{laurencon2024idefics3}, and Gemma~\cite{gemma2025} bring
$\leq$2B-parameter multimodal models to consumer hardware. LoRA~\cite{hu2022lora} and
QLoRA~\cite{dettmers2023qlora} allow task-specific adaptation with negligible inference
overhead---important when a scanner must not perceptibly slow down the host device.

\textbf{Calibration and silent failures.} Calibration of neural classifiers in the unimodal
setting is well studied~\cite{guo2017calibration}. We show that the dominant failure of
zero-shot small VLMs is not miscalibration but \emph{prediction collapse}---the model
functions as a constant-output device, not a miscalibrated one---and that standard
thresholded metrics (F1, accuracy) hide this from system evaluators. Zero-shot
vision--language transfer is established for recognition tasks~\cite{radford2021clip}; we
examine cross-\emph{category} transfer between two distinct harmful-content tasks.

%------------------------------------------------------------------------------
\section{Experimental Setup}\label{sec:setup}
%------------------------------------------------------------------------------
\textbf{Models.} SmolVLM-256M-Instruct ($\sim$489 MB weights), SmolVLM-500M-Instruct
($\sim$968 MB), and Gemma-4-E2B ($\sim$9,736 MB, $\sim$2B effective parameters). Memory
footprints are reported as peak allocated GPU VRAM at inference, which determines whether
the model can coexist with a host application on a device.

\textbf{Content categories.} \emph{Hate:} Hateful Memes~\cite{kiela2020hateful} balanced
dev split, $n{=}500$, AUROC primary metric. \emph{Sarcasm/irony:} MMSD2.0
\texttt{mmsd-v2}~\cite{qin2023mmsd2} test split, $n{=}500$. Both represent multimodal
content where the image and text must be jointly interpreted---exactly the class that
hash-based tools and text-only classifiers miss.

\textbf{Scanning protocol.} Each model receives a task instruction, the image, and a binary
question (``Is this content hateful/sarcastic?''). We extract $P(\text{yes})$ from the
softmax over the yes/no answer tokens---a continuous risk score for AUROC---and threshold
at 0.5 for hard flag/no-flag decisions. Gemma ships without a chat template and is queried
with a flattened raw prompt.

\textbf{Policy adaptation (LoRA).} We simulate per-organisation policy customisation:
LoRA~\cite{hu2022lora} rank 16 on all attention and MLP projections (1.85\% of parameters),
generative SFT on 3{,}000 examples for 2 epochs in bf16. At inference, the adapter adds
negligible overhead ($<$2\% parameter increase; identical latency profile). The adapter
could be shipped as a small policy file alongside the base model.

\textbf{Metrics.} AUROC (rank-based, threshold-free, robust to the label-prior collapse
documented in \S\ref{sec:collapse}); accuracy; and \emph{prediction-positive rate}
(\texttt{pred\_pos})---the fraction of inputs flagged as harmful, which exposes a broken
scanner. \textbf{Hardware.} All runs on a single NVIDIA A100 (80 GB); latency is the
median of 50 timed single-sample batch-1 runs after 5 warmup steps, representing a
single-item real-time scan query.

%------------------------------------------------------------------------------
\section{Finding 1: Memory Budget, Not Parameter Count, Governs Viability}\label{sec:scale}
%------------------------------------------------------------------------------
A practical on-device scanner must coexist with an operating system, other processes, and
potentially a user's applications on hardware that may have 4--16 GB of GPU-accessible
VRAM. Table~\ref{tab:zs} shows the zero-shot detection quality of each model against its
actual VRAM footprint. The frontier is \emph{non-monotone}: SmolVLM-256M sits at chance,
SmolVLM-500M is the best model on both content categories, and the larger Gemma-4-E2B
falls back toward chance while consuming $\sim$7$\times$ the memory and higher per-query
latency.

SmolVLM-500M therefore \emph{Pareto-dominates} Gemma-4-E2B as a scanner: it detects better
at a fraction of the cost. Critically, SmolVLM-500M's $\sim$1.5 GB peak VRAM footprint fits
on any device with a 4 GB GPU or even quantised to INT8 on CPU-only hardware, whereas
Gemma-4-E2B requires $\sim$10 GB, ruling out most consumer endpoints. At this scale,
architecture, instruction-tuning quality, and training data dominate raw parameter count.

\begin{table}[t]
\centering
\caption{Zero-shot detection quality vs.\ on-device cost (fp16, $n{=}500$, A100). AUROC is
the reliable metric (\S\ref{sec:collapse}); \texttt{pred\_pos} flags a broken scanner.
SmolVLM-500M fits on any 4 GB device; Gemma-4-E2B requires $\sim$10 GB.}
\label{tab:zs}
\begin{tabular}{lccccc}
\toprule
& \multicolumn{2}{c}{AUROC} & & Peak VRAM & Scan lat. \\
\cmidrule(lr){2-3}
Model & Hate & Sarcasm & \texttt{pred\_pos} & (MB) & (ms/item) \\
\midrule
SmolVLM-256M & 0.513 & 0.502 & 0.86--0.98 & \(\sim\)960  & 70 \\
SmolVLM-500M & \textbf{0.581} & \textbf{0.600} & 0.93--0.98 & \(\sim\)1{,}450 & 69 \\
Gemma-4-E2B  & 0.528 & 0.527 & 0.00        & \(\sim\)10{,}170 & 86 \\
\bottomrule
\end{tabular}
\end{table}

% TODO Fig.1: AUROC (y) vs. peak VRAM (x, log scale) scatter; Pareto frontier dashed;
% star = SmolVLM-500M + LoRA. Generate from results/leaderboard.jsonl.
\begin{figure}[t]
\centering
\fbox{\parbox{0.9\columnwidth}{\centering\vspace{1.0cm}\textit{Figure 1 (placeholder):
Detection quality vs.\ VRAM Pareto scatter (log-x); SmolVLM-500M dominates Gemma-4-E2B;
+LoRA point highest; 4 GB device boundary marked.}\vspace{1.0cm}}}
\caption{On-device viability: detection quality vs.\ memory footprint.}
\label{fig:pareto}
\end{figure}

%------------------------------------------------------------------------------
\section{Finding 2: Zero-Shot Models Fail Silently as Scanners}\label{sec:collapse}
%------------------------------------------------------------------------------
The AUROC values in Table~\ref{tab:zs} (0.50--0.60) already suggest weak separability, but
the \texttt{pred\_pos} column reveals a more dangerous failure: SmolVLM flags 86--98\% of
\emph{all} inputs as harmful; Gemma flags \emph{nothing} (0\%). Neither model behaves as a
classifier. They emit a near-constant output regardless of content.

This is not a miscalibrated scanner---it is a \emph{broken} one. A device running
zero-shot SmolVLM would generate near-universal false alarms, while one running Gemma would
silently pass every item. The failure is especially dangerous because \emph{standard metrics
hide it}. On a positive-skewed slice ($n{=}50$ items, majority harmful), SmolVLM-256M
yields an apparent \textbf{F1 of 0.89}---the kind of number that would make a security
engineer confident in the system. The same configuration on a balanced $n{=}500$ slice
yields AUROC 0.50 and sub-baseline accuracy. The strong F1 was an artifact of a
``flag-everything'' model meeting a mostly-harmful evaluation set.

The implication for security system evaluation is direct: \textbf{never validate a
content scanner on an imbalanced slice}. AUROC (threshold-free, balance-insensitive) and
\texttt{pred\_pos} (or the full score histogram, Fig.~\ref{fig:hist}) are the minimum
honest standard. A scanner that outputs the same score for benign and harmful content
provides no security guarantee regardless of how its F1 looks on a biased test set.

% TODO Fig.2: P(yes) score histograms per model -- SmolVLM mass near 1.0, Gemma mass near
% 0.0, showing the constant-output failure mode.
\begin{figure}[t]
\centering
\fbox{\parbox{0.9\columnwidth}{\centering\vspace{1.0cm}\textit{Figure 2 (placeholder):
$P(\text{yes})$ score histograms per model. SmolVLM: mass piled at 1.0 (flags everything).
Gemma: mass piled at 0.0 (flags nothing). A functioning scanner would show bimodal
separation.}\vspace{1.0cm}}}
\caption{Broken-scanner diagnostic: output-distribution collapse.}
\label{fig:hist}
\end{figure}

%------------------------------------------------------------------------------
\section{Finding 3: A Cheap Policy Adapter Fixes the Scanner---But Gains Differ 7$\times$
by Content Category}\label{sec:adapt}
%------------------------------------------------------------------------------
We fine-tune SmolVLM-500M with a LoRA adapter on each content category under an identical
recipe (\S\ref{sec:setup}). This simulates deploying a per-organisation policy file---a
small set of weights that customises the base model without changing it. Two effects stand
out (Table~\ref{tab:matrix}, diagonal).

\noindent\textbf{The adapter fixes the broken scanner first.} On both tasks, \texttt{pred\_pos}
moves from $\sim$0.95 toward the ground-truth prevalence rate ($0.36$ for hate, $0.29$ for
sarcasm). Fine-tuning's first-order effect is to stop the constant-flag collapse and
produce a model that actually \emph{abstains} on benign content---the minimum requirement
for a deployable scanner.

\noindent\textbf{Then the gains diverge by 7$\times$.} After de-biasing, sarcasm detection
improves dramatically ($0.600\to0.881$ AUROC, $+0.28$), while hate detection barely moves
($0.581\to0.619$, $+0.04$). Importantly, training loss converges to $\sim$0.1 for
\emph{both} tasks---the model memorises hate training examples as readily as sarcasm
ones---yet the hate adapter fails to generalise to held-out items. This is the signature of
the adversarial ``benign confounder'' design of the Hateful Memes
benchmark~\cite{kiela2020hateful}: surface-learnable train-set patterns do not transfer to
the evaluation set, which pairs each hateful item with a visually similar benign one. For
a practical scanner, this means that hate-meme detection is a fundamentally harder
engineering target than sarcasm: more diverse labelled data, harder negatives, or
confounder-aware training are needed, not simply more fine-tuning steps.

%------------------------------------------------------------------------------
\section{Finding 4: A Hate Adapter Bootstraps Sarcasm; Not the Reverse}\label{sec:transfer}
%------------------------------------------------------------------------------
A common deployment scenario: an organisation has labelled data for one content policy
(e.g.\ hate speech) and wants to extend coverage to a second category (e.g.\ sarcastic
harassment) with minimal additional labelling. Table~\ref{tab:matrix} (off-diagonal) shows
whether LoRA adapters transfer across categories.

The hate-trained adapter \emph{raises sarcasm detection} from the zero-shot floor of 0.600
to \textbf{0.668}---exceeding both the sarcasm zero-shot baseline and the hate adapter's
own in-domain hate score (0.619). The sarcasm-trained adapter, by contrast, leaves hate
detection essentially at the floor ($0.581\to0.594$). Transfer is strongly asymmetric.

The mechanism appears to be that hate fine-tuning instils a general
``read image and text jointly, commit to a decision'' behaviour required to resolve
confounder pairs---a skill that benefits any task requiring multimodal grounding, including
sarcasm. Sarcasm fine-tuning, in contrast, appears to latch onto a narrower
image--text incongruity cue that does not transfer back.

\textbf{Practical deployment corollary:} if building a scanner for a new content
category with limited labelled data, starting from a hate-detection adapter is a better
warm start than starting from a sarcasm one---or from zero-shot.

\begin{table}[t]
\centering
\caption{Cross-category transfer matrix, SmolVLM-500M (fp16, $n{=}500$): AUROC
(\texttt{pred\_pos}). Diagonal = in-domain (bold); off-diagonal = cross-category transfer.
Hate adapter cross-transfers positively to sarcasm; reverse does not hold.}
\label{tab:matrix}
\begin{tabular}{lcc}
\toprule
\textbf{Adapter trained on} $\downarrow$ & Hate AUROC & Sarcasm AUROC \\
\midrule
none (zero-shot floor) & 0.581 (0.93) & 0.600 (0.98) \\
Hate memes             & \textbf{0.619 (0.36)} & 0.668 (0.61) \\
Sarcasm                & 0.594 (0.65) & \textbf{0.881 (0.29)} \\
\bottomrule
\end{tabular}
\end{table}

%------------------------------------------------------------------------------
\section{Deployment Recommendations}\label{sec:recs}
%------------------------------------------------------------------------------
Our findings translate into concrete guidance for practitioners building on-device harmful
content scanners.

\noindent\textbf{Model selection.} Choose SmolVLM-500M (or similarly efficient $\sim$0.5B
instruction-tuned VLMs) over larger models. At $\sim$1.5 GB peak VRAM, it fits on any
consumer GPU at fp16 and is a candidate for quantised CPU-only deployment. Larger
($\sim$2B$+$) models provide no quality improvement at this scale and exceed the memory
budgets of most endpoint devices.

\noindent\textbf{Never ship zero-shot.} A zero-shot small VLM is a broken scanner. It
produces near-constant outputs that give a false sense of security and will not survive
a balanced audit. Always apply at minimum a small LoRA adapter trained on task-representative
data before deployment.

\noindent\textbf{Evaluate with AUROC and pred\_pos.} Do not validate a scanner on
skewed test slices or with F1/accuracy alone. Report AUROC (threshold-insensitive) and the
prediction-positive rate (\texttt{pred\_pos}); a functioning scanner must show
discrimination, not constant output.

\noindent\textbf{Budget for hate to be hard.} Adapters for hate-meme detection converge in
training but do not generalise well. Expect to invest in more diverse training data,
hard-negative mining, or confounder-aware augmentation. A sarcasm adapter converging to a
similar training loss is not an equivalent outcome.

\noindent\textbf{Bootstrap new categories from hate.} When labelled data for a new content
category is scarce, initialise from a hate-detection adapter rather than from zero-shot or
a sarcasm adapter.

%------------------------------------------------------------------------------
\section{Limitations}
%------------------------------------------------------------------------------
We study three models, two content categories, and single-seed $n{=}500$ balanced slices
(AUROC CI $\approx\pm0.03$--$0.04$). We benchmark on a server-grade A100; actual on-device
inference on laptop or mobile silicon will differ in latency and may require INT8/INT4
quantisation, which we have not yet evaluated. Gemma-4-E2B is evaluated without a chat
template (raw prompt), so its collapse may partly reflect prompt sensitivity rather than a
fundamental model property. We do not evaluate on truly novel or adversarial harmful
content beyond the two benchmarks, and we do not study detection evasion (adversarial
images). The hate-generalisation gap is inferred from train-loss/AUROC mismatch; a
full attribution study (e.g.\ confounder-stratified error analysis) would strengthen the
claim.

%------------------------------------------------------------------------------
\section{Conclusion}
%------------------------------------------------------------------------------
Small VLMs are a viable architectural basis for on-device, offline harmful-content
scanners---but only after lightweight policy adaptation. Zero-shot models silently fail as
scanners, and standard metrics conceal this failure. A single LoRA adapter ($<$2\% extra
parameters) converts a broken scanner into a functional one, though the detection gain
afterwards is 7$\times$ larger for sarcastic content than for adversarially designed hate
memes---a gap that reflects the relative difficulty of these content categories for compact
models, not a failure of the adaptation method. Transfer across categories is real but
asymmetric: a hate adapter is a useful bootstrap for sarcasm coverage when labelled data
is scarce, while the reverse does not hold. We release a reproducible harness and hope
these findings improve how on-device multimodal moderation tools are built, evaluated,
and reported.

\bibliographystyle{ACM-Reference-Format}
\bibliography{refs}

%==============================================================================
% TRIMMING TO 4 PAGES (if the CFP enforces it):
%  - Merge Findings 3 & 4 under one "Adaptation and Transfer" section; keep Table 2 only.
%  - Cut Figure 2 (keep the broken-scanner F1 paragraph as text).
%  - Compress Related Work to 4-5 sentences.
%  - Fold Deployment Recommendations into the Conclusion.
%  - Drop the per-contribution itemize in the intro.
%==============================================================================
\end{document}
